\section{The TRACE Framework: Dependence-Aware Evidence Gating}
\label{sec:framework}

TRACE sits between the detectors and the task manager. Detectors emit
per-frame reports. TRACE orders them, groups detections that refer to the
same physical object, and decides whether a task may close. Four pieces do
this work: a reconciler, a spatial association step, a decision policy, and an
audit record.

\subsection{Stream reconciliation with arrival watermarks}
Reports from different cameras arrive late and out of order. If the decision
depended on arrival order, the same evidence could close a task on one run and
hold it on the next.

The reconciler buffers each report and tracks the newest observation time it
has seen. The watermark is that time minus a fixed slack $\Delta t$. A report
is released once the watermark passes its observation time. Released reports
leave in one canonical order: observation time, then source identifier, then
report identifier. Any arrival permutation whose skew stays within $\Delta t$
therefore yields the same sequence, the same decisions, and the same trace
digest, byte for byte. A report that arrives behind the watermark cannot be
placed without rewriting decisions already released. TRACE stores it in a
separate late list. It is neither dropped nor merged into the ordered stream.
The slack is a parameter of the run. The sweep reported in
Section~\ref{sec:results} used $\Delta t = 1$\,s.

\subsection{Ground-plane spatial association}
Two cameras report the same vehicle at unrelated pixel coordinates, so TRACE
matches detections in the world. For each detection it casts a ray through the
bottom-centre pixel of the box, using the camera pose and the known camera
height, and intersects that ray with the ground plane. The result is a ground
point.

Matching is greedy. TRACE lists every pair of detections from different
agents that share a class label, sorts the pairs by ground distance, and
merges them in that order while the distance stays under 4\,m. A merge is
allowed only if the two clusters contain no agent in common, so one track
holds at most one detection per agent. The 4\,m limit absorbs the offset
between a vehicle's bottom edge as seen from opposite sides.

Some detections have no ground point. If the ray points at or above the
horizon, it never meets the ground. This happens for distant structures and
for false positives. TRACE keeps each such detection as its own single-source
track. It never matches them, so an unprojectable detection cannot be counted
as confirmed by a second camera.

\subsection{Dependence-aware lifecycle policy}
A vote count treats two agreeing reports as two units of evidence. The
calibration in Section~\ref{sec:results} shows that is wrong: the
agents miss together more often than independence allows. \texttt{MissRiskPolicy}
instead prices the chance that both agents missed something.

For an agent pair with detectors $d_A$ and $d_B$, the calibrated estimator
supplies the fitted ratio $\hat{\ratio}(d_A,d_B)$ and the two marginal miss
rates $p_A$ and $p_B$. The joint miss risk is
\begin{equation}
  P(\text{both miss}) = \hat{\ratio}(d_A,d_B)\; p_A\; p_B .
\end{equation}
The policy applies three actions.
\begin{itemize}
  \item \textsc{Wait}: fewer than two agents have reported, so the task waits
        for the second camera.
  \item \textsc{Verify}: any track has only one agent behind it, or the joint
        miss risk exceeds the tolerance $\tau$. The second case triggers a
        targeted secondary audit.
  \item \textsc{Close}: every track is co-observed and the joint miss risk is
        at most $\tau$.
\end{itemize}
Single-source tracks are checked first, because a lone detection is not
confirmed whatever the risk estimate says. With the independence estimator the
risk is zero, so the rule reduces to the single-source test alone. The
evaluation in Section~\ref{sec:results} ran in that mode, with
$\tau = 0.05$. The same policy code becomes dependence-aware when it is handed
the calibrated estimator. Evaluating that configuration is the open
experiment in Section~\ref{sec:limits}.

The policy handles two agents per task. With more, it reads the first two
source identifiers.

\subsection{Causal audit lineage}
Every decision records the identifiers of the reports that caused it, the
confirmed track identifiers, the reason code, the policy version, and the
estimator version. The estimator version is a fingerprint, the first 12 hex
characters of a SHA-256 over the frozen calibration numbers, so a decision
names exactly the estimator that produced it. The trace digest is the SHA-256
of the decisions serialised as canonical JSON (sorted keys, no whitespace).
Given the same reports, policy, and calibration, a replay reproduces the
digest. A changed digest means a changed input or a changed rule.
